\subsection*{Subsequences}

\begin{exposition}
Subsequences record what remains after passing to an infinite ordered selection
of indices. They are the basic tool for extracting structure from a sequence:
limits, oscillation, boundedness, compactness, and failure of convergence are
all detected by suitable subsequences.
\end{exposition}

% ---------------------------------------------------------
% Definition --- Strictly Increasing Index Map
% ---------------------------------------------------------

\begin{definitionbox}{Definition (Strictly Increasing Index Map)}
\begin{definition}[Strictly Increasing Index Map]
\label{def:strictly-increasing-index-map}
Let $\sigma:\mathbb{N}\to\mathbb{N}$ be a function. The function $\sigma$
is a \textbf{strictly increasing index map} if
\[
  k<\ell \Longrightarrow \sigma(k)<\sigma(\ell)
  \qquad
  \text{for all } k,\ell\in\mathbb{N}.
\]
\end{definition}
\end{definitionbox}

\begin{remark*}[Standard quantified statement]
\[
\sigma:\mathbb{N}\to\mathbb{N}\text{ is strictly increasing}
\Longleftrightarrow
\forall k,\ell\in\mathbb{N}\;(k<\ell\Longrightarrow \sigma(k)<\sigma(\ell)).
\]
\end{remark*}

\begin{remark*}[Definition predicate reading]
\[
\operatorname{StrictlyIncreasing}(\sigma)
\coloneqq
\forall k,\ell\in\mathbb{N}\;(k<\ell\Longrightarrow \sigma(k)<\sigma(\ell)).
\]
\end{remark*}

\begin{remark*}[Negated quantified statement]
\[
\sigma\text{ is not strictly increasing}
\Longleftrightarrow
\exists k,\ell\in\mathbb{N}\;(k<\ell\land \sigma(\ell)\le \sigma(k)).
\]
\end{remark*}

\begin{remark*}[Negation predicate reading]
\[
\neg\operatorname{StrictlyIncreasing}(\sigma)
\Longleftrightarrow
\exists k,\ell\in\mathbb{N}\;(k<\ell\land \sigma(\ell)\le \sigma(k)).
\]
\end{remark*}

\begin{remark*}[Interpretation]
The index map selects terms without repetition and without changing their
order. This is the structural difference between a subsequence and an
arbitrary rearrangement.
\end{remark*}

\begin{dependencies}
\begin{itemize}
  \item Functions
  \item Natural numbers
\end{itemize}
\end{dependencies}

% ---------------------------------------------------------
% Definition --- Subsequence
% ---------------------------------------------------------

\begin{definitionbox}{Definition (Subsequence)}
\begin{definition}[Subsequence]
\label{def:subsequence}
Let $(x_n)$ be a real sequence. A \textbf{subsequence} of $(x_n)$ is a
sequence $(x_{\sigma(k)})_{k\in\mathbb{N}}$, where
$\sigma:\mathbb{N}\to\mathbb{N}$ is a strictly increasing index map.
\end{definition}
\end{definitionbox}

\begin{remark*}[Standard quantified statement]
\[
\begin{aligned}
&(y_k)\text{ is a subsequence of }(x_n)\\
&\quad\Longleftrightarrow\quad
\exists \sigma:\mathbb{N}\to\mathbb{N}\;
\Bigl(
  \forall k,\ell\in\mathbb{N}\;(k<\ell\Longrightarrow \sigma(k)<\sigma(\ell))
  \land
  \forall k\in\mathbb{N}\;(y_k=x_{\sigma(k)})
\Bigr).
\end{aligned}
\]
\end{remark*}

\begin{remark*}[Definition predicate reading]
\[
\operatorname{Subsequence}(y_k,x_n)
\coloneqq
\exists \sigma\;
\bigl(
  \operatorname{StrictlyIncreasing}(\sigma)
  \land
  \forall k\in\mathbb{N}\;(y_k=x_{\sigma(k)})
\bigr).
\]
\end{remark*}

\begin{remark*}[Negated quantified statement]
\[
\begin{aligned}
&(y_k)\text{ is not a subsequence of }(x_n)\\
&\quad\Longleftrightarrow\quad
\forall \sigma:\mathbb{N}\to\mathbb{N}\;
\Bigl(
  \exists k,\ell\in\mathbb{N}\;(k<\ell\land \sigma(\ell)\le\sigma(k))
  \lor
  \exists k\in\mathbb{N}\;(y_k\neq x_{\sigma(k)})
\Bigr).
\end{aligned}
\]
\end{remark*}

\begin{remark*}[Negation predicate reading]
\[
\neg\operatorname{Subsequence}(y_k,x_n).
\]
\end{remark*}

\begin{remark*}[Interpretation]
A subsequence keeps the original chronological order but may skip terms. The
new index $k$ counts the selected terms, while $\sigma(k)$ records where each
selected term came from in the original sequence.
\end{remark*}

\begin{dependencies}
\begin{itemize}
  \item \hyperref[def:sequence]{Sequence}
  \item \hyperref[def:strictly-increasing-index-map]{Strictly Increasing Index Map}
\end{itemize}
\end{dependencies}

% ---------------------------------------------------------
% Definition --- Subsequential Limit
% ---------------------------------------------------------

\begin{definitionbox}{Definition (Subsequential Limit)}
\begin{definition}[Subsequential Limit]
\label{def:subsequential-limit}
Let $(x_n)$ be a real sequence and let $L\in\mathbb{R}$. The real number
$L$ is a \textbf{subsequential limit} of $(x_n)$ if there exists a
subsequence $(x_{\sigma(k)})$ such that
\[
  \lim_{k\to\infty} x_{\sigma(k)}=L.
\]
\end{definition}
\end{definitionbox}

\begin{remark*}[Standard quantified statement]
\[
\begin{aligned}
&L\text{ is a subsequential limit of }(x_n)\\
&\quad\Longleftrightarrow\quad
\exists \sigma:\mathbb{N}\to\mathbb{N}\;
\Bigl(
  \forall k,\ell\in\mathbb{N}\;(k<\ell\Longrightarrow \sigma(k)<\sigma(\ell))
  \land
  \lim_{k\to\infty}x_{\sigma(k)}=L
\Bigr).
\end{aligned}
\]
\end{remark*}

\begin{remark*}[Definition predicate reading]
\[
\operatorname{SubsequentialLimit}(x_n,L)
\coloneqq
\exists \sigma\;
\bigl(
  \operatorname{StrictlyIncreasing}(\sigma)
  \land
  \operatorname{ConvergentSequence}(x_{\sigma(k)},L)
\bigr).
\]
\end{remark*}

\begin{remark*}[Interpretation]
Subsequential limits detect limiting behavior along an infinite selection of
terms. A sequence may have no ordinary limit but still have one or more
subsequential limits.
\end{remark*}

\begin{dependencies}
\begin{itemize}
  \item \hyperref[def:subsequence]{Subsequence}
  \item \hyperref[def:convergent-sequence]{Convergence of a Sequence}
\end{itemize}
\end{dependencies}

% ---------------------------------------------------------
% Definition --- Has Convergent Subsequence
% ---------------------------------------------------------

\begin{definitionbox}{Definition (Has Convergent Subsequence)}
\begin{definition}[Has Convergent Subsequence]
\label{def:has-convergent-subsequence}
Let $(x_n)$ be a real sequence. The sequence $(x_n)$ \textbf{has a
convergent subsequence} if there exist a strictly increasing index map
$\sigma:\mathbb{N}\to\mathbb{N}$ and a real number $L$ such that
\[
  \lim_{k\to\infty}x_{\sigma(k)}=L.
\]
\end{definition}
\end{definitionbox}

\begin{remark*}[Standard quantified statement]
\[
\begin{aligned}
&(x_n)\text{ has a convergent subsequence}\\
&\quad\Longleftrightarrow\quad
\exists \sigma:\mathbb{N}\to\mathbb{N}\;\exists L\in\mathbb{R}\;
\Bigl(
  \forall k,\ell\in\mathbb{N}\;(k<\ell\Longrightarrow \sigma(k)<\sigma(\ell))
  \land
  \lim_{k\to\infty}x_{\sigma(k)}=L
\Bigr).
\end{aligned}
\]
\end{remark*}

\begin{remark*}[Definition predicate reading]
\[
\operatorname{HasConvergentSubsequence}(x_n)
\coloneqq
\exists L\in\mathbb{R}\;\operatorname{SubsequentialLimit}(x_n,L).
\]
\end{remark*}

\begin{remark*}[Interpretation]
This property asks for at least one convergent infinite selection. It does not
assert convergence of the whole sequence and does not require uniqueness of
the subsequential limit.
\end{remark*}

\begin{dependencies}
\begin{itemize}
  \item \hyperref[def:subsequential-limit]{Subsequential Limit}
\end{itemize}
\end{dependencies}

% ---------------------------------------------------------
% Theorem --- Subsequence Indices Dominate the Identity
% ---------------------------------------------------------

\begin{theorembox}{Theorem (Subsequence Indices Dominate the Identity)}
\begin{theorem}[Subsequence Indices Dominate the Identity]
\label{thm:subsequence-indices-dominate-identity}
Let $\sigma:\mathbb{N}\to\mathbb{N}$ be a strictly increasing index map.
Then
\[
  k\le \sigma(k)
  \qquad
  \text{for every } k\in\mathbb{N}.
\]

\smallskip
\noindent
\hyperref[prf:subsequence-indices-dominate-identity]{\textit{Go to proof.}}
\end{theorem}
\end{theorembox}

\begin{remark*}[Standard quantified statement]
\[
\forall \sigma:\mathbb{N}\to\mathbb{N}\;
\Bigl(
  \bigl[\forall k,\ell\in\mathbb{N}\;(k<\ell\Longrightarrow \sigma(k)<\sigma(\ell))\bigr]
  \Longrightarrow
  \bigl[\forall k\in\mathbb{N}\;(k\le\sigma(k))\bigr]
\Bigr).
\]
\end{remark*}

\begin{remark*}[Theorem predicate reading]
\[
\operatorname{StrictlyIncreasing}(\sigma)
\Longrightarrow
\forall k\in\mathbb{N}\;(k\le\sigma(k)).
\]
\end{remark*}

\begin{remark*}[Interpretation]
A subsequence cannot outrun the original indexing backwards. By the time the
subsequence has selected its $k$th term, it must have reached at least the
$k$th original position.
\end{remark*}

\begin{dependencies}
\begin{itemize}
  \item \hyperref[def:strictly-increasing-index-map]{Strictly Increasing Index Map}
  \item Well-ordering of $\mathbb{N}$
\end{itemize}
\end{dependencies}

% ---------------------------------------------------------
% Theorem --- Subsequences Preserve Limits
% ---------------------------------------------------------

\begin{theorembox}{Theorem (Subsequences Preserve Limits)}
\begin{theorem}[Subsequences Preserve Limits]
\label{thm:subsequences-preserve-limits}
Let $(x_n)$ be a real sequence and let $L\in\mathbb{R}$. If $x_n\to L$,
then every subsequence of $(x_n)$ converges to $L$.

\smallskip
\noindent
\hyperref[prf:subsequences-preserve-limits]{\textit{Go to proof.}}
\end{theorem}
\end{theorembox}

\begin{remark*}[Standard quantified statement]
\[
\begin{aligned}
&\forall (x_n)\subseteq\mathbb{R}\;\forall L\in\mathbb{R}\;
\forall \sigma:\mathbb{N}\to\mathbb{N}\\
&\quad\Bigl(
  x_n\to L
  \land
  \forall k,\ell\in\mathbb{N}\;(k<\ell\Longrightarrow\sigma(k)<\sigma(\ell))
\Bigr)
\Longrightarrow
\lim_{k\to\infty}x_{\sigma(k)}=L.
\end{aligned}
\]
\end{remark*}

\begin{remark*}[Theorem predicate reading]
\[
\operatorname{ConvergentSequence}(x_n,L)
\land
\operatorname{StrictlyIncreasing}(\sigma)
\Longrightarrow
\operatorname{ConvergentSequence}(x_{\sigma(k)},L).
\]
\end{remark*}

\begin{remark*}[Interpretation]
Subsequences cannot escape a limit already possessed by the whole sequence.
The domination of subsequence indices by the identity transfers every tail
estimate for $(x_n)$ to a tail estimate for $(x_{\sigma(k)})$.
\end{remark*}

\begin{dependencies}
\begin{itemize}
  \item \hyperref[def:subsequence]{Subsequence}
  \item \hyperref[def:convergent-sequence]{Convergence of a Sequence}
  \item \hyperref[thm:subsequence-indices-dominate-identity]{Subsequence Indices Dominate the Identity}
\end{itemize}
\end{dependencies}

% ---------------------------------------------------------
% Theorem --- Subsequential Limit of a Convergent Sequence
% ---------------------------------------------------------

\begin{theorembox}{Theorem (Subsequential Limit of a Convergent Sequence)}
\begin{theorem}[Subsequential Limit of a Convergent Sequence]
\label{thm:subsequential-limit-of-convergent-sequence}
Let $(x_n)$ be a real sequence and let $L\in\mathbb{R}$. If $x_n\to L$,
then $L$ is the only possible subsequential limit of $(x_n)$.

\smallskip
\noindent
\hyperref[prf:subsequential-limit-of-convergent-sequence]{\textit{Go to proof.}}
\end{theorem}
\end{theorembox}

\begin{remark*}[Standard quantified statement]
\[
\begin{aligned}
&\forall (x_n)\subseteq\mathbb{R}\;\forall L,K\in\mathbb{R}\\
&\quad\Bigl(
  x_n\to L
  \land
  K\text{ is a subsequential limit of }(x_n)
\Bigr)
\Longrightarrow
K=L.
\end{aligned}
\]
\end{remark*}

\begin{remark*}[Theorem predicate reading]
\[
\operatorname{ConvergentSequence}(x_n,L)
\land
\operatorname{SubsequentialLimit}(x_n,K)
\Longrightarrow
K=L.
\]
\end{remark*}

\begin{remark*}[Interpretation]
A convergent sequence has no hidden alternative limiting behavior. Every
convergent subsequence inherits the same limit, and uniqueness of limits then
forces equality.
\end{remark*}

\begin{dependencies}
\begin{itemize}
  \item \hyperref[thm:subsequences-preserve-limits]{Subsequences Preserve Limits}
  \item \hyperref[thm:uniqueness-of-limits]{Uniqueness of Limits}
\end{itemize}
\end{dependencies}

% ---------------------------------------------------------
% Theorem --- Divergence by Two Subsequential Limits
% ---------------------------------------------------------

\begin{theorembox}{Theorem (Divergence by Two Subsequential Limits)}
\begin{theorem}[Divergence by Two Subsequential Limits]
\label{thm:divergence-by-two-subsequential-limits}
Let $(x_n)$ be a real sequence. If $(x_n)$ has two distinct subsequential
limits, then $(x_n)$ does not converge.

\smallskip
\noindent
\hyperref[prf:divergence-by-two-subsequential-limits]{\textit{Go to proof.}}
\end{theorem}
\end{theorembox}

\begin{remark*}[Standard quantified statement]
\[
\begin{aligned}
&\forall (x_n)\subseteq\mathbb{R}\;\forall L,K\in\mathbb{R}\\
&\quad\Bigl(
  L\text{ is a subsequential limit of }(x_n)
  \land
  K\text{ is a subsequential limit of }(x_n)
  \land
  L\neq K
\Bigr)\\
&\quad\Longrightarrow\quad
\text{$(x_n)$ does not converge}.
\end{aligned}
\]
\end{remark*}

\begin{remark*}[Theorem predicate reading]
\[
\operatorname{SubsequentialLimit}(x_n,L)
\land
\operatorname{SubsequentialLimit}(x_n,K)
\land
L\neq K
\Longrightarrow
\neg\exists A\in\mathbb{R}\;\operatorname{ConvergentSequence}(x_n,A).
\]
\end{remark*}

\begin{remark*}[Interpretation]
Two different subsequential limits reveal persistent oscillation or splitting.
If the whole sequence had a limit, every subsequence would have to inherit
that same limit.
\end{remark*}

\begin{dependencies}
\begin{itemize}
  \item \hyperref[def:subsequential-limit]{Subsequential Limit}
  \item \hyperref[thm:subsequential-limit-of-convergent-sequence]{Subsequential Limit of a Convergent Sequence}
\end{itemize}
\end{dependencies}

% ---------------------------------------------------------
% Theorem --- Boundedness Passes to Subsequences
% ---------------------------------------------------------

\begin{theorembox}{Theorem (Boundedness Passes to Subsequences)}
\begin{theorem}[Boundedness Passes to Subsequences]
\label{thm:boundedness-passes-to-subsequences}
Every subsequence of a bounded real sequence is bounded.

\smallskip
\noindent
\hyperref[prf:boundedness-passes-to-subsequences]{\textit{Go to proof.}}
\end{theorem}
\end{theorembox}

\begin{remark*}[Standard quantified statement]
\[
\begin{aligned}
&\forall (x_n)\subseteq\mathbb{R}\;\forall \sigma:\mathbb{N}\to\mathbb{N}\\
&\quad\Bigl(
  \exists M>0\;\forall n\in\mathbb{N}\;(|x_n|\le M)
  \land
  \forall k,\ell\in\mathbb{N}\;(k<\ell\Longrightarrow\sigma(k)<\sigma(\ell))
\Bigr)\\
&\quad\Longrightarrow\quad
\exists M>0\;\forall k\in\mathbb{N}\;(|x_{\sigma(k)}|\le M).
\end{aligned}
\]
\end{remark*}

\begin{remark*}[Theorem predicate reading]
\[
\operatorname{BoundedSequence}(x_n)
\land
\operatorname{StrictlyIncreasing}(\sigma)
\Longrightarrow
\operatorname{BoundedSequence}(x_{\sigma(k)}).
\]
\end{remark*}

\begin{remark*}[Interpretation]
A subsequence cannot create new values; it only selects from the original
sequence. Any global bound for the original sequence is automatically a bound
for the selection.
\end{remark*}

\begin{dependencies}
\begin{itemize}
  \item \hyperref[def:bounded-sequence]{Bounded Sequence}
  \item \hyperref[def:subsequence]{Subsequence}
\end{itemize}
\end{dependencies}

% ---------------------------------------------------------
% Theorem --- Monotonicity Passes to Subsequences
% ---------------------------------------------------------

\begin{theorembox}{Theorem (Monotonicity Passes to Subsequences)}
\begin{theorem}[Monotonicity Passes to Subsequences]
\label{thm:monotonicity-passes-to-subsequences}
Every subsequence of an increasing real sequence is increasing, and every
subsequence of a decreasing real sequence is decreasing.

\smallskip
\noindent
\hyperref[prf:monotonicity-passes-to-subsequences]{\textit{Go to proof.}}
\end{theorem}
\end{theorembox}

\begin{remark*}[Standard quantified statement]
\[
\begin{aligned}
&\forall (x_n)\subseteq\mathbb{R}\;\forall \sigma:\mathbb{N}\to\mathbb{N}\\
&\quad\Bigl(
  \bigl[\forall n\in\mathbb{N}\;(x_n\le x_{n+1})\bigr]
  \land
  \bigl[\forall k,\ell\in\mathbb{N}\;(k<\ell\Longrightarrow\sigma(k)<\sigma(\ell))\bigr]
\Bigr)\\
&\quad\Longrightarrow\quad
\forall k\in\mathbb{N}\;(x_{\sigma(k)}\le x_{\sigma(k+1)}).
\end{aligned}
\]
\end{remark*}

\begin{remark*}[Theorem predicate reading]
\[
\operatorname{MonotoneIncreasingSequence}(x_n)
\land
\operatorname{StrictlyIncreasing}(\sigma)
\Longrightarrow
\operatorname{MonotoneIncreasingSequence}(x_{\sigma(k)}).
\]
\end{remark*}

\begin{remark*}[Interpretation]
Selecting terms in the original order preserves the order pattern already
present in the sequence. This is why the index map must be strictly increasing.
\end{remark*}

\begin{dependencies}
\begin{itemize}
  \item \hyperref[def:increasing-sequence]{Increasing Sequence}
  \item \hyperref[def:decreasing-sequence]{Decreasing Sequence}
  \item \hyperref[def:subsequence]{Subsequence}
\end{itemize}
\end{dependencies}

% ---------------------------------------------------------
% Theorem --- Subsequence of a Subsequence
% ---------------------------------------------------------

\begin{theorembox}{Theorem (Subsequence of a Subsequence)}
\begin{theorem}[Subsequence of a Subsequence]
\label{thm:subsequence-of-subsequence}
Every subsequence of a subsequence of $(x_n)$ is a subsequence of $(x_n)$.

\smallskip
\noindent
\hyperref[prf:subsequence-of-subsequence]{\textit{Go to proof.}}
\end{theorem}
\end{theorembox}

\begin{remark*}[Standard quantified statement]
\[
\begin{aligned}
&\forall (x_n)\subseteq\mathbb{R}\;\forall \sigma,\tau:\mathbb{N}\to\mathbb{N}\\
&\quad\Bigl(
  \bigl[\forall k,\ell\in\mathbb{N}\;(k<\ell\Longrightarrow\sigma(k)<\sigma(\ell))\bigr]
  \land
  \bigl[\forall k,\ell\in\mathbb{N}\;(k<\ell\Longrightarrow\tau(k)<\tau(\ell))\bigr]
\Bigr)\\
&\quad\Longrightarrow\quad
(x_{\sigma(\tau(k))})_{k\in\mathbb{N}}\text{ is a subsequence of }(x_n).
\end{aligned}
\]
\end{remark*}

\begin{remark*}[Theorem predicate reading]
\[
\operatorname{StrictlyIncreasing}(\sigma)
\land
\operatorname{StrictlyIncreasing}(\tau)
\Longrightarrow
\operatorname{Subsequence}(x_{\sigma(\tau(k))},x_n).
\]
\end{remark*}

\begin{remark*}[Interpretation]
Passing to subsequences is closed under iteration. The composed index map
$\sigma\circ\tau$ is still strictly increasing.
\end{remark*}

\begin{dependencies}
\begin{itemize}
  \item \hyperref[def:subsequence]{Subsequence}
  \item Composition of functions
\end{itemize}
\end{dependencies}

% ---------------------------------------------------------
% Theorem --- Eventual Properties Pass to Subsequences
% ---------------------------------------------------------

\begin{theorembox}{Theorem (Eventual Properties Pass to Subsequences)}
\begin{theorem}[Eventual Properties Pass to Subsequences]
\label{thm:eventual-properties-pass-to-subsequences}
Let $P(n)$ be a property of indices. If there exists $N\in\mathbb{N}$ such
that $P(n)$ holds for every $n\ge N$, then for every strictly increasing
index map $\sigma$ there exists $K\in\mathbb{N}$ such that $P(\sigma(k))$
holds for every $k\ge K$.

\smallskip
\noindent
\hyperref[prf:eventual-properties-pass-to-subsequences]{\textit{Go to proof.}}
\end{theorem}
\end{theorembox}

\begin{remark*}[Standard quantified statement]
\[
\begin{aligned}
&\Bigl(
  \exists N\in\mathbb{N}\;\forall n\ge N\;P(n)
  \land
  \forall k,\ell\in\mathbb{N}\;(k<\ell\Longrightarrow\sigma(k)<\sigma(\ell))
\Bigr)\\
&\quad\Longrightarrow\quad
\exists K\in\mathbb{N}\;\forall k\ge K\;P(\sigma(k)).
\end{aligned}
\]
\end{remark*}

\begin{remark*}[Theorem predicate reading]
\[
\operatorname{Eventually}(x_n,P)
\land
\operatorname{StrictlyIncreasing}(\sigma)
\Longrightarrow
\operatorname{Eventually}(x_{\sigma(k)},P).
\]
\end{remark*}

\begin{remark*}[Interpretation]
Subsequences eventually enter every tail of the original sequence. Any
property that holds on a tail of the original therefore holds on a tail of
the subsequence.
\end{remark*}

\begin{dependencies}
\begin{itemize}
  \item \hyperref[thm:subsequence-indices-dominate-identity]{Subsequence Indices Dominate the Identity}
  \item \hyperref[def:subsequence]{Subsequence}
\end{itemize}
\end{dependencies}

% ---------------------------------------------------------
% Theorem --- Frequent Properties Yield Subsequences
% ---------------------------------------------------------

\begin{theorembox}{Theorem (Frequent Properties Yield Subsequences)}
\begin{theorem}[Frequent Properties Yield Subsequences]
\label{thm:frequent-properties-yield-subsequences}
Let $P(n)$ be a property of indices. If for every $N\in\mathbb{N}$ there
exists $n\ge N$ such that $P(n)$ holds, then there exists a strictly
increasing index map $\sigma:\mathbb{N}\to\mathbb{N}$ such that
$P(\sigma(k))$ holds for every $k\in\mathbb{N}$.

\smallskip
\noindent
\hyperref[prf:frequent-properties-yield-subsequences]{\textit{Go to proof.}}
\end{theorem}
\end{theorembox}

\begin{remark*}[Standard quantified statement]
\[
\Bigl[
\forall N\in\mathbb{N}\;\exists n\ge N\;P(n)
\Bigr]
\Longrightarrow
\exists \sigma:\mathbb{N}\to\mathbb{N}\;
\Bigl(
  \forall k,\ell\in\mathbb{N}\;(k<\ell\Longrightarrow\sigma(k)<\sigma(\ell))
  \land
  \forall k\in\mathbb{N}\;P(\sigma(k))
\Bigr).
\]
\end{remark*}

\begin{remark*}[Theorem predicate reading]
\[
\operatorname{Frequently}(x_n,P)
\Longrightarrow
\exists \sigma\;\bigl(\operatorname{StrictlyIncreasing}(\sigma)
\land
\forall k\in\mathbb{N}\;P(\sigma(k))\bigr).
\]
\end{remark*}

\begin{remark*}[Interpretation]
An infinitely recurring property can be threaded into a subsequence. The
construction chooses one later witness at each step.
\end{remark*}

\begin{dependencies}
\begin{itemize}
  \item Well-ordering of $\mathbb{N}$
  \item \hyperref[def:subsequence]{Subsequence}
\end{itemize}
\end{dependencies}

% ---------------------------------------------------------
% Theorem --- Subsequential Limits Respect Bounds
% ---------------------------------------------------------

\begin{theorembox}{Theorem (Subsequential Limits Respect Bounds)}
\begin{theorem}[Subsequential Limits Respect Bounds]
\label{thm:subsequential-limits-respect-bounds}
Let $(x_n)$ be a real sequence and let $L$ be a subsequential limit of
$(x_n)$. If $m\le x_n\le M$ for every $n\in\mathbb{N}$, then
$m\le L\le M$.

\smallskip
\noindent
\hyperref[prf:subsequential-limits-respect-bounds]{\textit{Go to proof.}}
\end{theorem}
\end{theorembox}

\begin{remark*}[Standard quantified statement]
\[
\begin{aligned}
&\forall (x_n)\subseteq\mathbb{R}\;\forall L,m,M\in\mathbb{R}\\
&\quad\Bigl(
  L\text{ is a subsequential limit of }(x_n)
  \land
  \forall n\in\mathbb{N}\;(m\le x_n\le M)
\Bigr)
\Longrightarrow
m\le L\le M.
\end{aligned}
\]
\end{remark*}

\begin{remark*}[Theorem predicate reading]
\[
\operatorname{SubsequentialLimit}(x_n,L)
\land
\forall n\in\mathbb{N}\;(m\le x_n\le M)
\Longrightarrow
m\le L\le M.
\]
\end{remark*}

\begin{remark*}[Interpretation]
Bounds pass to subsequences, and limits preserve eventual inequalities. Hence
subsequential limits remain inside any closed interval that contains all terms.
\end{remark*}

\begin{dependencies}
\begin{itemize}
  \item \hyperref[thm:boundedness-passes-to-subsequences]{Boundedness Passes to Subsequences}
  \item \hyperref[thm:constant-comparison-for-sequence-limits]{Constant Comparison for Sequence Limits}
\end{itemize}
\end{dependencies}

% ---------------------------------------------------------
% Theorem --- Squeeze Passes to Subsequences
% ---------------------------------------------------------

\begin{theorembox}{Theorem (Squeeze Passes to Subsequences)}
\begin{theorem}[Squeeze Passes to Subsequences]
\label{thm:squeeze-passes-to-subsequences}
Let $(a_n)$, $(x_n)$, and $(b_n)$ be real sequences. If
$a_n\le x_n\le b_n$ eventually, then for every subsequence
$(x_{\sigma(k)})$ there exists $K\in\mathbb{N}$ such that
\[
  a_{\sigma(k)}\le x_{\sigma(k)}\le b_{\sigma(k)}
  \qquad
  \text{for every } k\ge K.
\]

\smallskip
\noindent
\hyperref[prf:squeeze-passes-to-subsequences]{\textit{Go to proof.}}
\end{theorem}
\end{theorembox}

\begin{remark*}[Standard quantified statement]
\[
\begin{aligned}
&\Bigl(
  \exists N\in\mathbb{N}\;\forall n\ge N\;(a_n\le x_n\le b_n)
  \land
  \forall k,\ell\in\mathbb{N}\;(k<\ell\Longrightarrow\sigma(k)<\sigma(\ell))
\Bigr)\\
&\quad\Longrightarrow\quad
\exists K\in\mathbb{N}\;\forall k\ge K\;
(a_{\sigma(k)}\le x_{\sigma(k)}\le b_{\sigma(k)}).
\end{aligned}
\]
\end{remark*}

\begin{remark*}[Theorem predicate reading]
\[
\operatorname{Eventually}(x_n,\;a_n\le x_n\le b_n)
\land
\operatorname{StrictlyIncreasing}(\sigma)
\Longrightarrow
\operatorname{Eventually}(x_{\sigma(k)},\;a_{\sigma(k)}\le x_{\sigma(k)}\le b_{\sigma(k)}).
\]
\end{remark*}

\begin{remark*}[Interpretation]
Squeeze hypotheses are tail hypotheses, so they survive passage to
subsequences.
\end{remark*}

\begin{dependencies}
\begin{itemize}
  \item \hyperref[thm:eventual-properties-pass-to-subsequences]{Eventual Properties Pass to Subsequences}
  \item \hyperref[thm:sequence-squeeze-theorem]{Sequence Squeeze Theorem}
\end{itemize}
\end{dependencies}

% ---------------------------------------------------------
% Theorem --- Monotone Subsequence Theorem
% ---------------------------------------------------------

\begin{theorembox}{Theorem (Monotone Subsequence Theorem)}
\begin{theorem}[Monotone Subsequence Theorem]
\label{thm:monotone-subsequence-theorem}
Every real sequence has a monotone subsequence.

\smallskip
\noindent
\hyperref[prf:monotone-subsequence-theorem]{\textit{Go to proof.}}
\end{theorem}
\end{theorembox}

\begin{remark*}[Standard quantified statement]
\[
\forall (x_n)\subseteq\mathbb{R}\;\exists \sigma:\mathbb{N}\to\mathbb{N}\;
\Bigl(
  \forall k,\ell\in\mathbb{N}\;(k<\ell\Longrightarrow\sigma(k)<\sigma(\ell))
  \land
  \bigl[
    \forall k\in\mathbb{N}\;(x_{\sigma(k)}\le x_{\sigma(k+1)})
    \lor
    \forall k\in\mathbb{N}\;(x_{\sigma(k+1)}\le x_{\sigma(k)})
  \bigr]
\Bigr).
\]
\end{remark*}

\begin{remark*}[Theorem predicate reading]
\[
\exists \sigma\;\bigl(
\operatorname{StrictlyIncreasing}(\sigma)
\land
\operatorname{MonotoneSequence}(x_{\sigma(k)})
\bigr).
\]
\end{remark*}

\begin{remark*}[Interpretation]
Every sequence contains an infinite ordered pattern. The proof uses the peak
dichotomy: either there are infinitely many terms larger than all later terms,
or there is an infinite rising chain.
\end{remark*}

\begin{dependencies}
\begin{itemize}
  \item \hyperref[def:monotone-sequence]{Monotone Sequence}
  \item \hyperref[def:subsequence]{Subsequence}
  \item Well-ordering of $\mathbb{N}$
\end{itemize}
\end{dependencies}

% ---------------------------------------------------------
% Theorem --- Bolzano-Weierstrass Theorem for Sequences
% ---------------------------------------------------------

\begin{theorembox}{Theorem (Bolzano-Weierstrass Theorem for Sequences)}
\begin{theorem}[Bolzano-Weierstrass Theorem for Sequences]
\label{thm:bolzano-weierstrass-sequences}
Every bounded real sequence has a convergent subsequence.

\smallskip
\noindent
\hyperref[prf:bolzano-weierstrass-sequences]{\textit{Go to proof.}}
\end{theorem}
\end{theorembox}

\begin{remark*}[Standard quantified statement]
\[
\forall (x_n)\subseteq\mathbb{R}\;
\Bigl(
  \exists M>0\;\forall n\in\mathbb{N}\;(|x_n|\le M)
  \Longrightarrow
  \exists \sigma:\mathbb{N}\to\mathbb{N}\;\exists L\in\mathbb{R}\;
  \bigl[
    \forall k,\ell\in\mathbb{N}\;(k<\ell\Longrightarrow\sigma(k)<\sigma(\ell))
    \land
    \lim_{k\to\infty}x_{\sigma(k)}=L
  \bigr]
\Bigr).
\]
\end{remark*}

\begin{remark*}[Theorem predicate reading]
\[
\operatorname{BoundedSequence}(x_n)
\Longrightarrow
\operatorname{HasConvergentSubsequence}(x_n).
\]
\end{remark*}

\begin{remark*}[Interpretation]
Bolzano-Weierstrass is the compactness principle for bounded real sequences.
Boundedness prevents escape, and the monotone subsequence theorem supplies an
ordered subsequence to which the Monotone Convergence Theorem applies.
\end{remark*}

\begin{dependencies}
\begin{itemize}
  \item \hyperref[def:bounded-sequence]{Bounded Sequence}
  \item \hyperref[thm:monotone-subsequence-theorem]{Monotone Subsequence Theorem}
  \item \hyperref[thm:boundedness-passes-to-subsequences]{Boundedness Passes to Subsequences}
  \item \hyperref[thm:bounded-monotone-sequence-equivalences]{Bounded Monotone Sequence Equivalences}
\end{itemize}
\end{dependencies}

% ---------------------------------------------------------
% Theorem --- Sequential Compactness of Closed Bounded Intervals
% ---------------------------------------------------------

\begin{theorembox}{Theorem (Sequential Compactness of Closed Bounded Intervals)}
\begin{theorem}[Sequential Compactness of Closed Bounded Intervals]
\label{thm:sequential-compactness-closed-bounded-interval}
Let $a,b\in\mathbb{R}$ with $a\le b$. Every sequence in $[a,b]$ has a
convergent subsequence whose limit belongs to $[a,b]$.

\smallskip
\noindent
\hyperref[prf:sequential-compactness-closed-bounded-interval]{\textit{Go to proof.}}
\end{theorem}
\end{theorembox}

\begin{remark*}[Standard quantified statement]
\[
\begin{aligned}
&\forall a,b\in\mathbb{R}\;\forall (x_n)\subseteq\mathbb{R}\\
&\quad\Bigl(
  a\le b
  \land
  \forall n\in\mathbb{N}\;(a\le x_n\le b)
\Bigr)\\
&\quad\Longrightarrow\quad
\exists \sigma:\mathbb{N}\to\mathbb{N}\;\exists L\in\mathbb{R}\;
\Bigl(
  \forall k,\ell\in\mathbb{N}\;(k<\ell\Longrightarrow\sigma(k)<\sigma(\ell))
  \land
  \lim_{k\to\infty}x_{\sigma(k)}=L
  \land
  a\le L\le b
\Bigr).
\end{aligned}
\]
\end{remark*}

\begin{remark*}[Theorem predicate reading]
\[
\bigl[\forall n\in\mathbb{N}\;(a\le x_n\le b)\bigr]
\Longrightarrow
\exists L\in[a,b]\;\operatorname{SubsequentialLimit}(x_n,L).
\]
\end{remark*}

\begin{remark*}[Interpretation]
Closed bounded intervals are sequentially compact: a sequence trapped in the
interval has a convergent subsequence, and closedness keeps the subsequential
limit inside the interval.
\end{remark*}

\begin{dependencies}
\begin{itemize}
  \item \hyperref[thm:bolzano-weierstrass-sequences]{Bolzano-Weierstrass Theorem for Sequences}
  \item \hyperref[thm:subsequential-limits-respect-bounds]{Subsequential Limits Respect Bounds}
\end{itemize}
\end{dependencies}